\section{Construction Pipeline Details}
\label{app:pipeline}

\begin{table}[t]
\centering
\small
\setlength{\tabcolsep}{4.5pt}
\renewcommand{\arraystretch}{1.50}
\begin{tabularx}{\linewidth}{@{}p{0.25\linewidth} X@{}}
\toprule
\textbf{Criterion} & \textbf{Requirement} \\
\midrule
Plausibility
& L1 and L2 options are plausible and distinct, with no strawmen or rewordings. \\
Re-ask
& Each layer adds a dimension absent from the initial question. \\
Counterfactuality
& L1 must eliminate or transform the core of L2, rather than leave it invariant. \\
\bottomrule
\end{tabularx}
\vspace{-2mm}
\caption{LLM judging rubric for chain injection. Both judges must pass three criteria for benchmark inclusion.}
\label{tab:chain-judge-rubric}
\vspace{-6mm}
\end{table}

This appendix expands the construction pipeline summarized in \S\ref{sec:construction}. For added clarity, We detail anchor selection, chain injection, drift injection, and evidence rewriting in order, specifying how candidates are selected, generated, filtered, and verified before their inclusion in \textsc{TIDE-Bench}.

\begin{table*}[t]
\centering
\small
\setlength{\tabcolsep}{5pt}
\renewcommand{\arraystretch}{1.60}
\begin{tabularx}{\linewidth}{@{}c p{0.15\linewidth} X p{0.22\linewidth}@{}}
\toprule
\textbf{Step} & \textbf{Check} & \textbf{Pass condition} & \textbf{Failure signal} \\
\midrule
1 & Output schema
& Required JSON fields are present.
& \texttt{json\_format} \\
2 & SQL parsing
& \texttt{post\_drift\_sql} parses and yields extractable commitments.
& \texttt{sql\_parse} \\
3 & SQL execution
& \texttt{post\_drift\_sql} executes on the target database.
& \texttt{sql\_execute} \\
4 & Non-empty result
& The post-drift query returns at least one row.
& \texttt{empty\_result} \\
5 & Observable effect
& Pre- and post-drift result sets differ.
& \texttt{identical\_result} \\
6 & Non-monotonicity
& The commitment diff satisfies $S_{\mathrm{pre}} \nsubseteq S_{\mathrm{post}}$.
& \texttt{no\_change}; \texttt{monotonic} \\
7 & Subtype match
& The realized drift subtype matches the assigned subtype.
& \texttt{subtype\_mismatch} \\
8 & Subtype detail
& For projection-level samples, the fine-grained subtype matches.
& \texttt{fine\_subtype} \\
\bottomrule
\end{tabularx}
\vspace{-0.5mm}
\caption{Sequential AST-level checks for validating drift candidates. Checks are applied in order; the final column reports the rejection signal emitted upon failure, and only candidates passing all checks undergo LLM judging.}
\label{tab:drift-ast}
\vspace{-4mm}
\end{table*}

\begin{table}[t]
\centering
\small
\setlength{\tabcolsep}{4.5pt}
\renewcommand{\arraystretch}{1.60}
\begin{tabular}{@{}cccccc@{}}
\toprule
Cond. & Evidence & Question & Dialogue & Acc. & $\Delta$\\
\midrule
\textbf{A} & orig. & orig.     & ---        & 51.2 & ---    \\
\textbf{B} & rew.  & ambiguous & \checkmark & 50.4 & $-0.8$ \\
\textbf{C} & rew.  & ambiguous & ---        & 31.2 & $-20.0$ \\
\bottomrule
\end{tabular}
\vspace{-0.5mm}
\caption{Reconstructability of rewritten evidence under three input settings ($n=695$). ``orig.''~/~``rew.'' denote original~/~rewritten BIRD evidence; \checkmark marks dialogue input; $\Delta$ is the overall accuracy change (pp) from A.}
\label{tab:evidence-ablation}
\vspace{-5mm}
\end{table}

\subsection{Chain Injection}
\label{app:pipeline-chain}
\noindent\textbf{Generation.}
We use GPT-5.1 to generate the chain candidates. For each BIRD entry, the model produces a JSON object with the initial question and the L1/L2 clarification pairs. The full prompt template and output schema are both shown in Figure~\ref{fig:chain-prompt}.

\namedpara{Judging.}
We score each candidate with two LLM judges, GPT-5.3-Codex and DeepSeek-V4-Pro. Each judge applies the three criteria in Table~\ref{tab:chain-judge-rubric}, in priority order. A candidate is accepted only if both judges pass all three criteria. The judges agree on 94.7\%, and 695/830 candidates pass (83.7\%).

\subsection{Drift Injection}
\label{app:pipeline-drift}

\noindent\textbf{Generation.}
We use GPT-5.1 to generate drift candidates. For each chain-validated anchor, the model receives a drift subtype: value-, constraint-, or projection-level. It returns a JSON object with the drift utterance and post-drift SQL. The prompt template and output schema are shown in Figure~\ref{fig:drift-prompt}.

\namedpara{AST-level validation.}
Each candidate is checked by a rule-based validator before LLM judging. The validator tests the operational definition of \S\ref{subsec:drift} and the scope constraints of the assigned subtype, as summarized in Table~\ref{tab:drift-ast}. The validation reuses the AST procedure from Appendix~\ref{app:commitment-extractor}; 554 of 695 validated chain candidates pass this stage (79.7\%).

\namedpara{Cross-validation.}
AST-passing candidates are scored by two LLM judges, Claude-Sonnet-4.6 and DeepSeek-V4-Pro. The judges receive the anchor SQL, drift utterance, and post-drift SQL. They apply two criteria in priority order: (i) \textbf{Drift fit}, requiring the post-drift SQL to match the assigned subtype and user drift; and (ii) \textbf{Stability}, requiring untouched units to remain unchanged. A candidate is accepted only if both judges pass both criteria, requiring agreement on validity. The judges agree on 96.6\%, and 514/554 candidates pass (92.8\%).

\subsection{Evidence Rewriting}
\label{app:pipeline-evidence}

\noindent\textbf{Motivation.}
BIRD's \texttt{evidence} field can reveal terms that should stay ambiguous in \textsc{TIDE-Bench}. For example, it may map ``top clients'' to ``clients ranked by spend,'' fixing the L1 choice ahead of dialogue. This early cue weakens the chain ambiguity we aim to test. The rewrite removes such cues while preserving schema-essential mappings.

\namedpara{Protocol.}
For each sample, GPT-5.1 rewrites the evidence using the source evidence, anchor SQL, and generated L1/L2 clarifications. The rewrite removes cues that resolve the ambiguity while retaining mappings needed for SQL interpretation.

\namedpara{Reconstructability.}
We test whether the rewritten evidence preserves schema-essential information without itself revealing the L1/L2 resolution. For 695 candidates, we evaluate Gemini-3-Flash under three settings in Table~\ref{tab:evidence-ablation}: the original question and evidence (A), the ambiguous question and rewritten evidence with dialogue (B), and the same inputs without dialogue (C). The small A--B gap of 0.8 pp shows that the rewritten evidence preserves what is needed for reconstruction when paired with dialogue. The large B--C gap of 19.2 pp shows that the rewritten evidence alone does not leak resolution.

